\subsection{Proof of Theorem \ref{thm:exact}}

By exact invariance, \(\mu(Qz_1,\ldots,Qz_n)\in\ran Q\). Since \(QQ^*\) is
the identity on \(\ran Q\),
\[
Q\bar\mu(z_1,\ldots,z_n)
=QQ^*\mu(Qz_1,\ldots,Qz_n)
=\mu(Qz_1,\ldots,Qz_n).
\]
For a partial composition, apply this identity first to the inner law. Its
lifted output is exactly the ambient inner output, so the outer law receives
the same arguments on both sides. Induction on the number of internal nodes
proves the statement for every finite ordered tree.

\subsection{Proof of the identity-inheritance corollary}

Every term \(T_a\) is a finite tree. The preceding induction gives
\[
QT_{a,\bar\mu}(z)=T_{a,\mu}(Qz).
\]
Linearity in the coefficients \(c_a\) gives the asserted identity for
\(F=\sum_ac_aT_a\). Since \(Q\) is injective, a zero ambient polynomial
restricts to a zero reduced polynomial.

\subsection{Proof of Theorem \ref{thm:approx}}

For a node \(v\), let \(F_v\) be the unprojected ambient value, \(R_v\) the
projected recursive value, \(k_v\) the number of internal nodes below \(v\),
and \(L_v\) the product of norms of leaves below \(v\). We prove
\[
\norm{F_v},\norm{R_v}\leq M^{k_v}L_v,\qquad
\norm{F_v-R_v}\leq k_v\rho M^{k_v-1}L_v.
\]
For a leaf, \(F_v=R_v=Qz\), \(k_v=0\), and the claim is immediate. At an
internal vertex,
\[
\begin{aligned}
\norm{F_v-R_v}
&=\norm{\mu(F_{v_1},\ldots,F_{v_n})
-P\mu(R_{v_1},\ldots,R_{v_n})}\\
&\leq \norm{\mu(F_{v_1},\ldots,F_{v_n})
-\mu(R_{v_1},\ldots,R_{v_n})}\\
&\quad+\norm{(I-P)\mu(R_{v_1},\ldots,R_{v_n})}.
\end{aligned}
\]
The first term is bounded by the multilinear telescoping inequality
\[
M\sum_i\norm{F_{v_i}-R_{v_i}}
\prod_{j\ne i}\max(\norm{F_{v_j}},\norm{R_{v_j}}).
\]
Every \(R_{v_i}\) belongs to \(\ran P\), so the second term is at most
\(\rho\prod_i\norm{R_{v_i}}\). Substitute the induction bounds. With
\(k_v=1+\sum_i k_{v_i}\) and \(L_v=\prod_iL_{v_i}\), the \(i\)-th telescoping
term is at most
\[
k_{v_i}\rho M^{k_v-1}L_v,
\]
and the local residual term is at most
\[
\rho M^{k_v-1}L_v.
\]
Summing gives \(k_v\rho M^{k_v-1}L_v\). The norm bounds follow from
\(\norm{\mu(u_1,\ldots,u_n)}\leq M\prod_i\norm{u_i}\) and contractivity of
\(P\). Finally, \(\norm{PF_T-R_T}\leq\norm{F_T-R_T}\).

\subsection{Proof of Corollary \ref{cor:assocbound}}

The left and right associator terms are two-node ternary trees. Apply
Theorem~\ref{thm:approx} to each term and use the triangle inequality:
\[
\norm{QA_{\bar\mu}-PA_\mu}
\leq 2M\rho\prod_{j=1}^5\norm{z_j}
+2M\rho\prod_{j=1}^5\norm{z_j}.
\]
This is the stated \(4M\rho\) bound.

\subsection{Spectral-snapping justification}

The selected eigenspaces of \(A\) and \(A+E\) are separated from the
threshold when \(\norm{E}_2\leq\gamma/2\), so their dimensions agree by
Weyl's eigenvalue perturbation bound. Applying a standard Davis--Kahan
sin-theta estimate with the threshold separation and replacing the exact
separation by the conservative lower bound \(\gamma/2\) yields a constant no
larger than \(4/\gamma\) in the stated regime. The theorem intentionally
records the safe threshold constant rather than a sharp variant.

\subsection{Counterexample verification}

For the no-invariance construction, the reduced inner output is zero, while
the ambient inner output on \((e_0,e_0)\) is \(e_1\), and the outer law maps
\((e_1,e_0)\) to \(e_0\). Hence the two compositions differ. For the
no-gap construction, the two diagonal eigenvalues cross the threshold after
an \(O(\delta)\) perturbation, and the snapped coordinate projectors are
\(\operatorname{diag}(0,1)\) and \(\operatorname{diag}(1,0)\), whose operator
distance is one.
